\subsection{韦尔-菲廷分解}

设 A 是环，M 是一个 A-模，u 是 M 的一个自同态。对任意整数 \(p \geqslant 0\) ，我们把 \(u^{p}\) 的核记作 \(N_{p}\) 。子模序列 \((\mathrm{N}_{p})\) 是递增的，其并是 M 的一个在 u 下稳定的子模 \(N_{\infty}\) 。对每个整数 \(p \geqslant 0\) ，我们有 \(\mathrm{N}_{p+1} = \overline{u}^{-1}(\mathrm{N}_{p})\) ，因而关系 \(N_{p} = N_{p+1}\) 蕴涵 \(N_{p+1} = N_{p+2}\) 。于是，或者序列 \((\mathrm{N}_{p})\) 严格递增，或者存在整数 \(p \geqslant 0\) ，使得 \(N_{0}, \ldots, N_{p}\) 互不相同且 \(N_{p} = N_{\infty}\) 。

对任意整数 \(q \geqslant 0\) ，我们把 \(u^{q}\) 的像记作 \(I_{q}\) 。子模序列 ( \(I_{q}\) ) 是递降的，其交是 M 的一个在 u 下稳定的子模 \(I_{\infty}\) 。对每个整数 \(q \geqslant 0\) ，我们有 \(u(\mathrm{I}_{q}) = \mathrm{I}_{q+1}\) ，因而关系 \(I_{q} = I_{q+1}\) 蕴涵 \(I_{q+1} = I_{q+2}\) 。于是，或者序列 ( \(I_{q}\) ) 严格递降，或者存在整数 \(q \geqslant 0\) ，使得 \(I_{0}, \ldots, I_{q}\) 互不相同且 \(I_{q} = I_{\infty}\) 。

命题 2. --- a) 设序列 \((\mathrm{N}_{p})\) 是平稳的。则我们有 \(N_{\infty} \cap I_{\infty} = 0\) ，u 在 \(I_{\infty}\) 上的限制是单射，且 u 诱导 \(N_{\infty}\) 上的一个幂零自同态。

\begin{enumerate}
\def\labelenumi{\alph{enumi})}
\setcounter{enumi}{1}
\item
  设序列 \((\mathrm{I}_p)\) 是平稳的。则我们有 \(\mathbf{M} = \mathrm{N}_{\infty} + \mathrm{I}_{\infty}\) 及 \(u(\mathrm{I}_{\infty}) = \mathrm{I}_{\infty}\) 。
\item
  （``韦尔--菲廷分解''——有时称为``菲廷分解''）设序列 \((\mathrm{N}_{p})\) 与 \((\mathrm{I}_{p})\) 都是平稳的。则 M 是在 u 下稳定的子模 \(N_{\infty}\) 与 \(I_{\infty}\) 的直和，且 u 诱导 \(N_{\infty}\) 上的一个幂零自同态以及 \(I_{\infty}\) 上的一个自同构。
\end{enumerate}

设 p 是使得 \(N_{p} = N_{\infty}\) 的自然数，令 \(v = u^{p}\) 。由构造，v 与 \(v^{2}\) 有相同的核 \(N_{\infty}\) ，而 \(I_{p}\) 是 v 的像。对 \(N_{\infty} \cap I_{p}\) 中的 x，存在 \(y \in M\) ，使得 \(x = v(y)\) ，从而 \(v^{2}(y) = v(x) = 0\) ；因此我们有 \(y \in N_{\infty}\) ，于是 x = 0。特别地，我们有 \(N_{\infty} \cap I_{\infty} = 0\) 。由于 u 的核 \(N_{1}\) 含于 \(N_{\infty}\) ，u 在 \(I_{\infty}\) 上的限制是单射；另一方面，我们有 \(u^{p}(N_{\infty}) = 0\) 。这就证明了 a)。

设 \(q\) 是使得 \(\mathrm{I}_q = \mathrm{I}_{\infty}\) 的自然数，令 \(w = u^q\) 。则 \(w\) 与 \(w^2\) 有相同的像 \(\mathrm{I}_{\infty}\) ，而 \(\mathrm{N}_q\) 是 \(w\) 的核。设 \(x \in \mathrm{M}\) 。我们有 \(w(x) \in \mathrm{I}_{\infty}\) ，故存在 \(y \in \mathrm{M}\) ，使得 \(w(x) = w^2(y)\) ；于是我们有 \(x - w(y) \in \mathrm{N}_q\) ，从而 \(\mathrm{M} = \mathrm{N}_q + \mathrm{I}_{\infty}\) ，更有 \(\mathrm{M} = \mathrm{N}_{\infty} + \mathrm{I}_{\infty}\) 。我们有 \(u(\mathrm{I}_{\infty}) = u(\mathrm{I}_q) = \mathrm{I}_{q+1} = \mathrm{I}_{\infty}\) 。这就证明了 b)。

断言 c) 由 a) 与 b) 立即得出。

注. --- 1) 设 p 是使得 \(N_{p} = N_{p+1}\) 的整数；上面给出的证明表明 \(N_{\infty} \cap I_{p} = 0\) ，且 u 在 \(I_{p}\) 上的限制是单射。同样，设 q 是使得 \(I_{q} = I_{q+1}\) 的整数；则我们有 \(N_{q} + I_{\infty} = M\) ，且由 u 通过过渡到商而得到的 \(M/N_{q}\) 的自同态是满射。

\begin{enumerate}
\def\labelenumi{\arabic{enumi})}
\setcounter{enumi}{1}
\tightlist
\item
  设 M 是两个在 u 下稳定的子模 N 与 I 的直和，且 u 诱导 N 上的一个幂零自同态 \(u_{N}\) 与 I 上的一个自同构。则我们有 \(N_{\infty} = N\) 与 \(I_{\infty} = I\) ，且序列 \((N_{p})\) 与 \((I_{p})\) 都是平稳的。此外，下列整数相等：
\end{enumerate}

\(\alpha)\) 最小的整数 \(p \geqslant 0\) ，它使 \(\mathrm{N}_p = \mathrm{N}_{\infty}\) ，

\(\beta)\) 最小的整数 \(q \geqslant 0\) ，它使 \(\mathrm{I}_q = \mathrm{I}_{\infty}\) ，

\(\gamma)\) 最小的整数 \(r \geqslant 0\) ，它使 \((u_{\mathrm{N}})^{r} = 0\) 。

\begin{enumerate}
\def\labelenumi{\arabic{enumi})}
\setcounter{enumi}{2}
\tightlist
\item
  若 A-模 M 是诺特的，则断言 a) 的假设被满足；若 M 是阿廷的，则断言 b) 的假设被满足；由 VIII, p.~2 的命题 1，若 M 有有限长度，则断言 c) 的假设被满足。
\end{enumerate}

\textbf{推论 1. --- 设 A 是环，M 是一个 A-模。}

\begin{enumerate}
\def\labelenumi{\alph{enumi})}
\item
  若模 M 是诺特的，则 M 的每个满射自同态都是双射的。
\item
  若模 M 是阿廷的，则 M 的每个单射自同态都是双射的。
\item
  若模 M 有有限长度，则 M 的每个单射或满射自同态都是双射的。
\item
  若环 A 交换且 A-模 M 是有限生成的，则 M 的每个满射自同态都是双射的。
\end{enumerate}

设 u 是 A-模 M 的一个自同态。我们使用本小节开头引入的记号。若自同态 u 是满射的，则我们有 \(I_{\infty} = M\) 。于是断言 a) 由命题 2 a) 与注 3 得出。同样，若自同态 u 是单射的，则我们有 \(N_{\infty} = 0\) 。于是断言 b) 由命题 2 b) 与注 3 得出。断言 c) 由 a) 与 b) 立即得出。

现在设环 A 交换、A-模 M 有限生成、且自同态 u 是满射的。我们来证明 u 是单射的。设 x 是 M 的一个满足 \(u(x) = 0\) 的元素。选取 A-模 M 的一个有限生成族 \((x_{i})_{i \in \mathrm{I}}\) ，并对每个 \(i \in I\) ，选取 M 的一个元素 \(y_{i}\) ，使得 \(u(y_{i}) = x_{i}\) 。存在 A 的元素的族 \((a_{i})_{i \in \mathrm{I}}, (b_{ij})_{(i,j) \in \mathrm{I} \times \mathrm{I}}\) 与 \((c_{ij})_{(i,j) \in \mathrm{I} \times \mathrm{I}}\) ，使得我们有

\[
x = \sum_ {i \in \mathrm{I}} a _ {i} x _ {i}, \quad y _ {j} = \sum_ {i \in \mathrm{I}} b _ {i j} x _ {i}, \quad u (x _ {j}) = \sum_ {i \in \mathrm{I}} c _ {i j} x _ {i}
\]

对每个 \(j \in I\) 。设 \(A'\) 是 A 的一个包含元素 \(a_i\) 、 \(b_{ij}\) 与 \(c_{ij}\) 的诺特子环（VIII, p.~12, 推论 3）。设 \(M'\) 是 M 的 \(A'\) -子模，由族 \((x_i)_{i \in I}\) 生成。我们有 \(u(x_j) \in M'\) 、 \(y_j \in M'\) 以及 \(u(y_j) = x_j\) （对每个 \(j \in I\) ）；于是 u 由限制给出一个满射自同态 \(u'\) （\(A'\) -模 \(M'\) 的）。由于环 \(A'\) 是诺特的，有限生成的 \(A'\) -模 \(M'\) 是诺特的（VIII, p.~7, 命题 4 a)）。由 a)，自同态 \(u'\) （\(M'\) 的）是双射的。由构造，x 属于 \(M'\) ，且我们有 \(u'(x) = u(x) = 0\) 。因此 x = 0，这就证明了 d)。

推论 2. --- 在一个左诺特环中，每个左可逆或右可逆的元素都是可逆的。

事实上，考虑元素 \(x\) 、 \(y\) （某个左诺特环 \(A\) 的），它们满足 \(xy = 1\) 。分别用 \(\pmb{\delta}(x)\) 与 \(\pmb{\delta}(y)\) 记自同态 \(a \mapsto ax\) 与 \(a \mapsto ay\) （A-模 \(A_s\) 的）。我们有 \(\pmb{\delta}(y) \circ \pmb{\delta}(x) = 1_{A_s}\) ；因此 \(\pmb{\delta}(y)\) 是满射的。由推论 1 a)，\(\pmb{\delta}(y)\) 是双射的。于是自同态 \(\pmb{\delta}(x)\) 是 \(\pmb{\delta}(y)\) 的逆双射，且我们有 \(yx = (\pmb{\delta}(x) \circ \pmb{\delta}(y))(1) = 1\) 。推论得证。
% semantic-fragments: legacy-input-seam
\relax{}
